\documentclass[a4paper,12pt]{article}

\usepackage[T2A]{fontenc}
\usepackage[utf8]{inputenc}
\usepackage[russian]{babel}
\usepackage{amsmath,amssymb,amsthm,mathtools}
\usepackage{helvet}
\renewcommand{\familydefault}{\sfdefault}
\usepackage{icomma}
\usepackage[a4paper,left=2.2cm,right=2.2cm,top=2.1cm,bottom=2.2cm]{geometry}
\usepackage{microtype}
\usepackage{setspace}
\onehalfspacing
\usepackage{parskip}
\usepackage{tikz}
\usetikzlibrary{calc}
\usepackage{tkz-euclide}
\usepackage{pgfplots}
\pgfplotsset{compat=1.18}
\usepackage{tabularx,booktabs,longtable}
\usepackage{enumitem}
\usepackage{xcolor}
\usepackage{fancyhdr}
\usepackage{hyperref}

\definecolor{accent}{HTML}{008080}
\definecolor{darkaccent}{HTML}{004D4D}
\definecolor{textgray}{HTML}{333333}
\definecolor{softgray}{HTML}{707070}
\definecolor{rulegray}{HTML}{D5DEDE}
\color{textgray}

\hypersetup{
    colorlinks=true,
    linkcolor=darkaccent,
    urlcolor=darkaccent,
    pdfauthor={Лёвин Артём Александрович},
    pdftitle={Чек-лист: квадратные корни и преобразование выражений}
}

\pagestyle{fancy}
\fancyhf{}
\fancyhead[L]{\small\color{softgray}Квадратные корни и преобразование выражений}
\fancyhead[R]{\small\color{softgray}ОГЭ}
\fancyfoot[C]{\small\color{softgray}\thepage}
\renewcommand{\headrulewidth}{0.4pt}
\renewcommand{\headrule}{\hbox to\headwidth{\color{rulegray}\leaders\hrule height \headrulewidth\hfill}}
\renewcommand{\footrulewidth}{0pt}

\setlist[itemize]{leftmargin=1.7em,itemsep=0.35em,topsep=0.35em}
\setlist[enumerate]{leftmargin=2em,itemsep=0.55em,topsep=0.4em}
\newlist{checklist}{itemize}{1}
\setlist[checklist]{label={\(\square\)},leftmargin=1.8em,itemsep=0.45em,topsep=0.4em}

\newcommand{\topic}[1]{%
    \vspace{0.7em}
    {\large\bfseries\color{darkaccent}#1}\par
    \vspace{0.25em}
    {\color{rulegray}\hrule}
    \vspace{0.65em}
}
\newcommand{\important}[1]{{\bfseries\color{darkaccent}#1}}

\begin{document}

\begin{titlepage}
\thispagestyle{empty}
\begin{center}
\vspace*{2.1cm}
{\Large\color{softgray}ЧЕК-ЛИСТ}
\vspace{0.7cm}

{\fontsize{25}{31}\selectfont\bfseries\color{darkaccent}
Квадратные корни\\[0.25em]
и преобразование выражений
}

\vspace{0.7cm}
{\large
Сравнение чисел, свойства корней,\\
формулы сокращённого умножения
}

\vspace{0.9cm}
{\color{accent}\rule{0.58\linewidth}{1pt}}
\vspace{1.1cm}

{\large Подготовка к ОГЭ по математике}
\vfill

{\large
Лёвин Артём Александрович\\[0.25em]
эксклюзивно для Ярослава Верещагина
}

\vspace{0.8cm}
{\large \today}
\vspace{1.4cm}

\begin{tikzpicture}[x=1cm,y=1cm]
    \draw[accent,line width=1.15pt]
    (-6,0)
    .. controls (-3.7,0.65) and (-2.2,-0.55) .. (0,0)
    .. controls (2.1,0.55) and (3.8,-0.45) .. (6,0.12);
\end{tikzpicture}
\end{center}
\end{titlepage}

\section*{\color{darkaccent}Основной чек-лист}

\topic{1. Квадратный корень: смысл и ограничения}

Для действительных чисел запись
\[
\sqrt{a}
\]
имеет смысл при
\[
a\ge 0.
\]

Квадратный корень из \(a\) --- это неотрицательное число, квадрат которого равен \(a\):
\[
\sqrt{a}=b
\quad\Longleftrightarrow\quad
b\ge0,\qquad b^2=a.
\]

Например,
\[
\sqrt{25}=5,
\qquad
\sqrt{0}=0.
\]

\begin{checklist}
    \item Подкоренное выражение проверено: оно должно быть неотрицательным.
    \item Помню, что арифметический квадратный корень всегда неотрицателен.
    \item Не заменяю \(\sqrt{a^2}\) на \(a\) без проверки знака \(a\).
\end{checklist}

\important{Ключевая формула:}
\[
\boxed{\sqrt{a^2}=|a|}
\]

Например,
\[
\sqrt{(-7)^2}=7.
\]

Если заранее известно, что \(a\ge0\), то
\[
\sqrt{a^2}=a.
\]

\topic{2. Основные свойства квадратных корней}

При \(a\ge0\), \(b\ge0\):
\[
\boxed{\sqrt a\cdot\sqrt b=\sqrt{ab}}
\]

В обратную сторону:
\[
\boxed{\sqrt{ab}=\sqrt a\cdot\sqrt b}
\]

При \(a\ge0\), \(b>0\):
\[
\boxed{\frac{\sqrt a}{\sqrt b}=\sqrt{\frac ab}}
\]

Особенно полезно:
\[
\boxed{\sqrt a\cdot\sqrt a=a},
\qquad a\ge0.
\]

Пример:
\[
\sqrt{11}\cdot\sqrt{11}=11.
\]

\important{Важно.}
Формулы произведения и частного квадратных корней в действительных числах применяются только при допустимых подкоренных выражениях.

\topic{3. Разложение подкоренного выражения}

Главная практическая идея:
\[
\sqrt{ab}=\sqrt a\sqrt b.
\]

Число под корнем удобно раскладывать так, чтобы появился полный квадрат.

Например:
\[
\sqrt{18}=\sqrt{9\cdot2}=3\sqrt2.
\]

Аналогично:
\[
\sqrt{20}=2\sqrt5,
\qquad
\sqrt{27}=3\sqrt3,
\qquad
\sqrt{45}=3\sqrt5.
\]

\begin{checklist}
    \item Ищу внутри подкоренного выражения множитель \(4,9,16,25,36,\ldots\).
    \item Выношу полный квадрат из-под корня.
    \item После преобразования проверяю, можно ли объединить одинаковые радикалы.
\end{checklist}

Пример:
\[
\sqrt2\bigl(\sqrt{18}+\sqrt2\bigr).
\]

Раскрываем скобки:
\[
\sqrt2\cdot\sqrt{18}+\sqrt2\cdot\sqrt2.
\]

Можно сразу объединить корни:
\[
\sqrt{36}+2=6+2=8.
\]

Можно получить тот же результат через \(\sqrt{18}=3\sqrt2\):
\[
\sqrt2\cdot3\sqrt2+2=6+2=8.
\]

\topic{4. Перегруппировка произведений с корнями}

В выражениях с несколькими корнями полезно переставлять множители так, чтобы одинаковые радикалы оказались рядом.

Например:
\[
5\sqrt{11}\cdot2\sqrt2\cdot\sqrt{22}.
\]

Так как
\[
\sqrt{22}=\sqrt2\sqrt{11},
\]
то
\[
5\sqrt{11}\cdot2\sqrt2\cdot\sqrt2\sqrt{11}.
\]

Группируем:
\[
5\cdot2\cdot(\sqrt2\sqrt2)\cdot(\sqrt{11}\sqrt{11}).
\]

Получаем
\[
10\cdot2\cdot11=220.
\]

Следовательно,
\[
\boxed{5\sqrt{11}\cdot2\sqrt2\cdot\sqrt{22}=220}.
\]

\important{Контроль арифметики.}
Числовые коэффициенты перед корнями также входят в произведение и не должны теряться.

\topic{5. Формулы сокращённого умножения}

При работе с корнями особенно полезны три формулы:
\[
\boxed{(a+b)^2=a^2+2ab+b^2},
\]
\[
\boxed{(a-b)^2=a^2-2ab+b^2},
\]
\[
\boxed{(a-b)(a+b)=a^2-b^2}.
\]

Вместо \(a\) и \(b\) могут стоять целые выражения, в том числе корни.

Например:
\[
(\sqrt{11}-\sqrt3)(\sqrt{11}+\sqrt3).
\]

Используем разность квадратов:
\[
(\sqrt{11})^2-(\sqrt3)^2=11-3=8.
\]

Ещё один пример:
\[
(\sqrt{11}+3)^2-6\sqrt{11}.
\]

Раскрываем квадрат суммы:
\[
11+6\sqrt{11}+9-6\sqrt{11}=20.
\]

Следовательно,
\[
\boxed{(\sqrt{11}+3)^2-6\sqrt{11}=20}.
\]

\begin{checklist}
    \item Если вместо \(a\) стоит \(\sqrt{11}\), рассматриваю весь радикал как единое выражение.
    \item При возведении суммы в квадрат не забываю среднее слагаемое \(2ab\).
    \item Перед раскрытием длинных скобок проверяю, нет ли формулы сокращённого умножения.
\end{checklist}

\topic{6. Квадрат выражения с корнем}

Если всё выражение возводится в квадрат, квадрат относится ко всему основанию.

Например:
\[
(4\sqrt3)^2=4^2(\sqrt3)^2=16\cdot3=48.
\]

Поэтому
\[
\frac{(4\sqrt3)^2}{60}=\frac{48}{60}=\frac45.
\]

\important{Типичная ошибка:}
\[
(4\sqrt3)^2\ne4\cdot3.
\]

\topic{7. Степень под квадратным корнем}

Для целого \(k\ge0\):
\[
\boxed{\sqrt{a^{2k}}=|a|^k}.
\]

Например:
\[
\sqrt{a^4}=a^2,
\qquad
\sqrt{a^8}=a^4.
\]

Если известно, что \(a\ge0\), удобно использовать запись
\[
\sqrt{a^n}=a^{n/2}.
\]

Например:
\[
\sqrt{6^4}=6^{4/2}=6^2=36.
\]

\important{Ограничение.}
Механическое правило «поделить степень на \(2\)» нельзя применять, игнорируя знак основания. Например,
\[
\sqrt{a^2}=|a|,
\]
а не просто \(a\).

\topic{8. Полезные приближённые значения}

Для быстрого оценивания на ОГЭ полезно помнить:
\[
\boxed{\sqrt2\approx1{,}41},
\qquad
\boxed{\sqrt3\approx1{,}73},
\qquad
\boxed{\sqrt5\approx2{,}24}.
\]

\begin{table}[ht]
\centering
\caption{Быстрые оценки распространённых корней}
\renewcommand{\arraystretch}{1.25}
\begin{tabular}{cc}
\toprule
Корень & Приближённое значение \\
\midrule
\(\sqrt2\) & \(1{,}41\) \\
\(\sqrt3\) & \(1{,}73\) \\
\(\sqrt5\) & \(2{,}24\) \\
\bottomrule
\end{tabular}
\end{table}

Приближения удобны для быстрой оценки. Если сравниваемые числа очень близки, следует либо увеличить точность приближения, либо перейти к точному сравнению.

\topic{9. Сравнение чисел разного вида}

Пусть требуется сравнить
\[
\frac5{12},\qquad0{,}44,\qquad\frac7{15},\qquad\frac{\sqrt2}{3}.
\]

Есть два основных пути.

\subsection*{\color{darkaccent}Путь 1. Десятичные приближения}
\[
\frac5{12}\approx0{,}4167,
\qquad
0{,}44=0{,}44,
\qquad
\frac7{15}=0{,}4666\ldots,
\]
\[
\frac{\sqrt2}{3}\approx\frac{1{,}41}{3}\approx0{,}47.
\]

Отсюда видно:
\[
\frac5{12}<0{,}44<\frac7{15}<\frac{\sqrt2}{3}.
\]

При слишком грубом приближении \(\sqrt2\approx1{,}4\) получилось бы
\[
\frac{\sqrt2}{3}\approx0{,}4666\ldots,
\]
и различить два последних числа было бы трудно.

\subsection*{\color{darkaccent}Путь 2. Точное сравнение}

Переведём десятичную дробь:
\[
0{,}44=\frac{44}{100}=\frac{11}{25}.
\]

Для знаменателей \(12,25,15,3\):
\[
\operatorname{НОК}(12,25,15,3)=300.
\]

Действительно,
\[
12=2^2\cdot3,
\qquad
25=5^2,
\qquad
15=3\cdot5,
\qquad
3=3,
\]
поэтому
\[
\operatorname{НОК}=2^2\cdot3\cdot5^2=300.
\]

Получаем:
\[
\frac5{12}=\frac{125}{300},
\qquad
\frac{11}{25}=\frac{132}{300},
\qquad
\frac7{15}=\frac{140}{300},
\qquad
\frac{\sqrt2}{3}=\frac{100\sqrt2}{300}.
\]

Остаётся точно сравнить \(140\) и \(100\sqrt2\). Оба числа положительны, поэтому сравниваем квадраты:
\[
140^2=19600,
\qquad
(100\sqrt2)^2=20000.
\]

Так как \(19600<20000\), то
\[
140<100\sqrt2.
\]

Следовательно,
\[
\boxed{\frac5{12}<0{,}44<\frac7{15}<\frac{\sqrt2}{3}}.
\]

\topic{10. Как точно сравнивать положительные выражения с корнями}

Если \(A\ge0\) и \(B\ge0\), то
\[
A<B\quad\Longleftrightarrow\quad A^2<B^2.
\]

Например, сравним
\[
\frac7{15}\quad\text{и}\quad\frac{\sqrt2}{3}.
\]

Умножаем обе положительные величины на \(15\):
\[
7\quad\text{и}\quad5\sqrt2.
\]

Возводим в квадрат:
\[
7^2=49,
\qquad
(5\sqrt2)^2=50.
\]

Следовательно,
\[
\boxed{\frac7{15}<\frac{\sqrt2}{3}}.
\]

\important{Важно.}
Сравнение квадратов эквивалентно сравнению самих чисел только после проверки их знаков.

\topic{11. Общий знаменатель и НОК}

Чтобы привести несколько дробей к общему знаменателю:
\begin{enumerate}
    \item разложите знаменатели на простые множители;
    \item возьмите каждый простой множитель в наибольшей встречающейся степени;
    \item перемножьте выбранные множители;
    \item найдите дополнительные множители для каждой дроби.
\end{enumerate}

Например:
\[
12=2^2\cdot3,
\qquad
15=3\cdot5,
\qquad
25=5^2.
\]

Поэтому
\[
\operatorname{НОК}(12,15,25)=2^2\cdot3\cdot5^2=300.
\]

\begin{checklist}
    \item При домножении знаменателя обязательно домножаю на то же число числитель.
    \item Перед поиском общего знаменателя сокращаю дроби, если это возможно.
    \item Для нескольких знаменателей удобнее искать НОК через простые множители.
\end{checklist}

\topic{12. Корректное сокращение дробей}

Сокращать можно только \important{множители}.

Например:
\[
\frac{16}{20}=\frac{4\cdot4}{5\cdot4}=\frac45.
\]

Нельзя сокращать отдельные части суммы:
\[
\frac{a+b}{a}
\]
не превращается в \(1+b\).

\topic{13. Итоговый алгоритм для выражений с корнями}

Перед вычислением сложного выражения:
\begin{enumerate}
    \item Проверьте область допустимых значений.
    \item Посмотрите, есть ли формула сокращённого умножения.
    \item Разложите подкоренные числа на удобные множители.
    \item Вынесите полные квадраты из-под корня.
    \item Если корни перемножаются, попробуйте объединить их.
    \item Сгруппируйте одинаковые корни.
    \item Выполните сокращение только по множителям.
    \item При сравнении близких чисел переходите к точному методу.
    \item В конце выполните арифметическую проверку.
\end{enumerate}

\topic{14. Типичные ошибки}

\begin{checklist}
    \item Не пишу \(\sqrt{a^2}=a\) без учёта знака \(a\).
    \item Не забываю возводить числовой коэффициент в квадрат: \((4\sqrt3)^2=16\cdot3\).
    \item Не теряю коэффициенты перед корнями при перегруппировке.
    \item Не использую грубое приближение, если числа получаются почти одинаковыми.
    \item Не сокращаю слагаемые через дробную черту.
    \item При приведении к общему знаменателю изменяю одновременно числитель и знаменатель.
    \item Перед сравнением квадратов проверяю, что сравниваемые величины неотрицательны.
    \item Помню: \(\sqrt{a+b}\ne\sqrt a+\sqrt b\) в общем случае.
\end{checklist}

Например:
\[
\sqrt{9+16}=5,
\qquad
\sqrt9+\sqrt{16}=3+4=7.
\]

\clearpage
\section*{\color{darkaccent}Тренировочный блок}

\textbf{Путь А.}
Выполните задания без калькулятора. Решения оформляйте последовательно, с указанием основных преобразований.

\subsection*{\color{darkaccent}Средний уровень}
\begin{enumerate}
    \item Найдите значение:
    \[\sqrt{12}\cdot\sqrt3.\]

    \item Найдите значение:
    \[3\sqrt5\cdot2\sqrt{20}.\]

    \item Найдите значение:
    \[(\sqrt{13}-\sqrt5)(\sqrt{13}+\sqrt5).\]
\end{enumerate}

\subsection*{\color{darkaccent}Выше среднего}
\begin{enumerate}[resume]
    \item Найдите значение:
    \[\sqrt3\bigl(\sqrt{27}+\sqrt3\bigr).\]

    \item Найдите значение:
    \[(\sqrt7+2)^2-4\sqrt7.\]

    \item Найдите значение:
    \[\frac{(5\sqrt2)^2}{50}.\]

    \item Упростите выражение:
    \[\frac{\sqrt{45}+\sqrt{20}}{\sqrt5}.\]

    \item Сравните без использования десятичных приближений:
    \[\frac7{15}\quad\text{и}\quad\frac{\sqrt2}{3}.\]

    \item Расположите числа в порядке возрастания:
    \[\frac5{12},\qquad0{,}44,\qquad\frac7{15},\qquad\frac{\sqrt2}{3}.\]
\end{enumerate}

\subsection*{\color{darkaccent}Сложный уровень}
\begin{enumerate}[resume]
    \item Найдите значение:
    \[
    \frac{(\sqrt{14}+\sqrt6)(\sqrt{14}-\sqrt6)}{\sqrt2(\sqrt8+\sqrt2)}
    +
    \frac{(\sqrt{11}+3)^2-6\sqrt{11}}{10}.
    \]
\end{enumerate}

\clearpage
\section*{\color{darkaccent}Ответы к тренировочному блоку}

\begin{table}[ht]
\centering
\caption{Ответы}
\renewcommand{\arraystretch}{1.4}
\begin{tabularx}{\linewidth}{>{\centering\arraybackslash}p{1.4cm}>{\centering\arraybackslash}X}
\toprule
№ & Ответ \\
\midrule
1 & \(6\) \\
2 & \(60\) \\
3 & \(8\) \\
4 & \(12\) \\
5 & \(11\) \\
6 & \(1\) \\
7 & \(5\) \\
8 & \(\displaystyle\frac7{15}<\frac{\sqrt2}{3}\) \\
9 & \(\displaystyle\frac5{12}<0{,}44<\frac7{15}<\frac{\sqrt2}{3}\) \\
10 & \(\displaystyle\frac{10}{3}\) \\
\bottomrule
\end{tabularx}
\end{table}

\vfill
{\small\color{softgray}
При проверке заданий с квадратными корнями отдельно контролируйте область допустимых значений, знаки выражений и корректность применения свойств радикалов.
}

\end{document}
